\documentclass[11pt,a4paper]{article}
\usepackage[T2A]{fontenc}
\usepackage[utf8]{inputenc}
\usepackage[russian,english]{babel}
\usepackage{amsmath,amssymb,amsthm,mathtools}
\usepackage{setspace}
\usepackage{enumitem}
\usepackage[hidelinks]{hyperref}
\onehalfspacing
\hyphenpenalty=10000\exhyphenpenalty=10000\emergencystretch=3em
\newcommand{\head}[1]{\par\bigskip\noindent\textbf{\underline{#1}}\quad}
\newcommand{\rudate}{14 сентября 2026 года}
\newcommand{\sheettitle}[3]{% title, subtitle, homework-flag (empty or "Homework")
  \begin{center}
  {\huge\bfseries #1\par}\vspace{1.2em}
  {\Large #2\par}\vspace{0.8em}
  \if\relax\detokenize{#3}\relax\else{\Large #3\par}\vspace{0.8em}\fi
  {\foreignlanguage{russian}{\rudate}\par}
  \end{center}\vspace{0.5em}}
\newcommand{\src}[1]{\unskip\nobreak\hfil\penalty50\hskip1em\hbox{}\nobreak\hfill(#1){\parfillskip=0pt\par}}
\newcommand{\fl}[1]{\left\lfloor #1\right\rfloor}
\DeclareMathOperator{\lcm}{lcm}
\DeclareMathOperator{\ord}{ord}
\begin{document}
\sheettitle{Number theory -- 2}{Orders and primitive roots}{Homework}

\head{Theoretical minimum.} For $\gcd(a,n)=1$ the order $\ord_n(a)$ is the least $d>0$ with
$a^d\equiv1$; then $a^m\equiv1$ if and only if the order divides $m$, and the order divides
$\varphi(n)$. Hence every prime divisor of $a^n-1$ either divides some $a^d-1$ with $d\mid n$, $d<n$,
or is $\equiv1\pmod n$. The group $(\mathbb Z/p\mathbb Z)^*$ is cyclic, because $x^d=1$ has at most
$d$ roots over the field $\mathbb F_p$ and $\sum_{d\mid p-1}\varphi(d)=p-1$; a primitive root modulo
$n$ exists exactly for $n=1,2,4,p^k,2p^k$, while
$(\mathbb Z/2^k\mathbb Z)^*\cong\mathbb Z/2\times\mathbb Z/2^{k-2}$ for $k\geqslant3$. Fixing a
primitive root turns multiplication into addition, so $a$ is an $m$-th power modulo $p$ if and only
if $a^{(p-1)/\gcd(m,p-1)}\equiv1$. Zsygmondy's theorem, which may be quoted, provides a primitive
prime divisor of $a^n-b^n$ with two exceptions.

\medskip\noindent\rule{\linewidth}{0.4pt}

\begin{enumerate}[leftmargin=2.2em,itemsep=0.6em]
\item Find all prime numbers $p$ for which there is a unique $a\in\{1,2,\dots,p\}$ with $a^3-3a+1$
divisible by $p$. \src{IMC}

\item Prove that for a prime $p>3$ the number of $k\in\{1,\dots,p-2\}$ with $I(k+1)<I(k)$ is greater
than $p/4-1$, where $I(k)\in\{1,\dots,p-1\}$ is the inverse of $k$ modulo $p$. \src{Putnam 2025}

\item Let $b_0=0$ and $b_{n+1}=2b_n^2+b_n+1$ for $n\geqslant0$. Prove that for each $k\geqslant1$ the
number $b_{2^{k+1}}-2b_{2^k}$ is divisible by $2^{2k+2}$ but not by $2^{2k+3}$. \src{Putnam 2025}

\item Let $k$ be a positive integer and $n=2^k+1$. Prove that $n$ is prime if and only if there
exist a permutation $a_1,\dots,a_{n-1}$ of $1,\dots,n-1$ and integers $g_1,\dots,g_{n-1}$ such that
$n$ divides $g_i^{a_i}-a_{i+1}$ for every $i$, where $a_n=a_1$. \src{IMO shortlist 2011, N8}

\item Determine the multiplicative order of $2$ modulo a Mersenne prime, prove that for
$k\geqslant2$ a prime divisor of $2^{2^k}+1$ is $\equiv1\pmod{2^{k+2}}$, and deduce that the Fermat
numbers are pairwise coprime; then give the resulting proof that there are infinitely many primes.

\item Let $g$ be a primitive root modulo an odd prime $p$. Prove that exactly $\varphi(p-1)$ residues
are primitive roots, that the product of all primitive roots is $\equiv1\pmod p$ for $p>3$, and that
the sum of all primitive roots is $\equiv\mu(p-1)\pmod p$.

\item Prove that for every $n$ there are infinitely many primes $p$ with $p\equiv1\pmod n$, by
considering the prime divisors of the values of the cyclotomic polynomial $\Phi_n(x)$ at integers,
and explain why this argument does not reach the residue classes other than $1$.

\item Determine all integers $m\geqslant2$ such that every $n$ with
$m/3\leqslant n\leqslant m/2$ divides the binomial coefficient $\binom{n}{m-2n}$.
\src{IMO shortlist 2012, N3}

\item Prove that there are infinitely many positive integers $n$ such that the largest prime
divisor of $n^4+n^2+1$ equals the largest prime divisor of $(n+1)^4+(n+1)^2+1$.
\src{IMO shortlist 2013, N3}

\item Prove that for every prime $p>100$ and every integer $r$ there exist integers $a$ and $b$
such that $p$ divides $a^2+b^5-r$. \src{IMO shortlist 2012, N8}
\end{enumerate}
\end{document}
